\chapter{共振与非绝热跃迁}
\label{ch:resonance-lz}

离散态若与连续谱耦合，原来的能级会获得位移和有限寿命；两个随时间靠近的能级若彼此耦合，系统又可能在避免交叉处改变绝热分支。前者是定态散射问题，后者是含时演化问题，却都需要区分被直接观测的态和其余状态。投影消元把 Hilbert 空间分成保留子空间 \(P\) 与其余子空间 \(Q\)，解出 \(Q\) 分量，再将它反馈到 \(P\) 中；定态反馈形成能量依赖的自能，含时反馈一般带有记忆。Fano--Feshbach 模型用于计算离散态与连续谱的共振，Landau--Zener 模型则在给定封闭二能级哈密顿量下求避免交叉的含时跃迁。
\section{Feshbach 投影}

设正交投影算符

\[
P^2=P,
\qquad
Q^2=Q,
\qquad
PQ=0,
\qquad
P+Q=\mathbb I.
\]

对复谱参数 $z$，先取 $z\in\rho(H)\cap\rho(H_{QQ})$。有限维情形可直接使用分块矩阵；无界算符还要求所选子空间与共同定义域相容，以下连续谱模型的非对角耦合为有界有限秩算符。完整预解式 $G(z)=(z-H)^{-1}$ 的 $P$ 块满足 Schur 补恒等式
\begin{equation}
PG(z)P
=\left[z-H_{PP}-H_{PQ}(z-H_{QQ})^{-1}H_{QP}\right]^{-1}_{P}.
\label{eq:feshbach-schur-resolvent}
\end{equation}
这里右端逆只在 $P$ 子空间中取。定态有效哈密顿量、共振极点、能级位移和寿命都只是这一预解式恒等式在不同解析区域与边界条件下的读法。

\begin{proof}
设 Hilbert 空间分解 \(\mathcal H=\mathcal H_P\oplus\mathcal H_Q\)，正交投影 \(P+Q=I\)。对 \(z\in\rho(H)\)，预解式 \(G(z)=(z-H)^{-1}\) 满足
\[
 (z-H)|\Psi\rangle=|\eta_P\rangle,
 \qquad Q|\eta_P\rangle=0.
\]
把 \(|\Psi\rangle\) 分块为 \(|\Psi\rangle=|\psi_P\rangle+|\psi_Q\rangle\)，其中 \(|\psi_P\rangle=P|\Psi\rangle\)、\(|\psi_Q\rangle=Q|\Psi\rangle\)。将方程分别左乘 \(P\) 和 \(Q\)，利用 \(PHP=H_{PP}\)、\(PHQ=H_{PQ}\) 等分块记号，得到耦合方程组
\begin{align*}
 (z-H_{PP})|\psi_P\rangle-H_{PQ}|\psi_Q\rangle
 &=|\eta_P\rangle,\\
 -H_{QP}|\psi_P\rangle+(z-H_{QQ})|\psi_Q\rangle
 &=0.
\end{align*}
第一式是 \(P\) 空间分量，第二式是 \(Q\) 空间分量，两式通过非对角块 \(H_{PQ}\)、\(H_{QP}\) 耦合。

只要 \(z\) 位于 \(H_{QQ}\) 的预解集，\((z-H_{QQ})\) 可逆，第二式精确解为
\[
 |\psi_Q\rangle
 =(z-H_{QQ})^{-1}H_{QP}|\psi_P\rangle.
\]


将 \(|\psi_Q\rangle\) 代回第一式，
\[
 (z-H_{PP})|\psi_P\rangle
 -H_{PQ}(z-H_{QQ})^{-1}H_{QP}|\psi_P\rangle
 =|\eta_P\rangle,
\]
整理为
\[
 \left[z-H_{PP}
 -H_{PQ}(z-H_{QQ})^{-1}H_{QP}\right]|\psi_P\rangle
 =|\eta_P\rangle.
\]
方括号为 $z-H_{\rm eff}(z)$，其中
\[
 H_{\rm eff}(z)=H_{PP}+H_{PQ}(z-H_{QQ})^{-1}H_{QP},
\]
因此 \(|\psi_P\rangle=(z-H_{\rm eff}(z))^{-1}|\eta_P\rangle\)。另一方面，由 \(|\eta_P\rangle\in\mathcal H_P\) 和 \(G(z)=(z-H)^{-1}\)，
\[
 |\psi_P\rangle=P|\Psi\rangle=PG(z)|\eta_P\rangle=PG(z)P|\eta_P\rangle.
\]
两个表达式对任意 \(|\eta_P\rangle\in\mathcal H_P\) 成立，故算子恒等
\[
 PG(z)P=(z-H_{\rm eff}(z))^{-1}
 =\left[z-H_{PP}-H_{PQ}(z-H_{QQ})^{-1}H_{QP}\right]^{-1},
\]
即\eqnrefs{eq:feshbach-schur-resolvent}。修正项 \(H_{PQ}(z-H_{QQ})^{-1}H_{QP}\) 描述从 \(P\) 跳到 \(Q\)、传播再跳回 \(P\) 的二阶虚过程。

边界条件：出射波 \(z\to E+\ii0^+\)、入射波 \(z\to E-\ii0^+\)。由 Sokhotski--Plemelj 公式，
\[
 (E\pm\ii0^+-H_{QQ})^{-1}
 =\operatorname{p.v.}\frac1{E-H_{QQ}}
 \mp\ii\pi\delta(E-H_{QQ}),
\]
取出射支，对归一化的保留态 $|\psi_P\rangle$，自能的虚部给出能量宽度
\[
 \Gamma_\psi(E)=-2\operatorname{Im}\langle\psi_P|\Sigma(E)|\psi_P\rangle
 =2\pi\langle\psi_P|H_{PQ}\delta(E-H_{QQ})H_{QP}|\psi_P\rangle\ge0.
\]
这里的边界值及 $\delta(E-H_{QQ})$ 先按谱测度分布理解：对能量作非负测试函数积分后，所得二次型非负。若耦合向量的谱测度在所讨论的区间有绝对连续密度，$\Gamma_\psi(E)$ 才是几乎处处定义的非负宽度函数；不能将式中的 $\delta$ 解释成在每个实能量处都有界的算符。只有在自能随能量变化缓慢、单个极点主导的条件下，才以 $E_R-\ii\Gamma/2$ 近似共振极点。其时间振幅为 $\ee^{-\ii E_Rt/\hbar}\ee^{-\Gamma t/(2\hbar)}$，概率衰变率为 $\Gamma/\hbar$；$\Gamma$ 本身具有能量量纲。
\end{proof}

定态 Schr\"odinger 方程

\[
(E-H)|\Psi\rangle=0
\]

分块为

\begin{align*}
(E-H_{PP})P|\Psi\rangle
&=H_{PQ}Q|\Psi\rangle,
\\
(E-H_{QQ})Q|\Psi\rangle
&=H_{QP}P|\Psi\rangle.
\end{align*}

散射问题要求出射边界条件，因此预解式应取

\[
G_Q^{(+)}(E)
\equiv(E+\ii0^+-H_{QQ})^{-1}.
\]

散射态还含入射的齐次解。设 $|\chi_E\rangle\in\mathcal H_Q$ 是给定入射通道的广义态，满足 $(E-H_{QQ})|\chi_E\rangle=0$。在它之上叠加出射响应，第二式解为

\[
Q|\Psi\rangle
=|\chi_E\rangle+G_Q^{(+)}(E)H_{QP}P|\Psi\rangle.
\]

代回 $P$ 空间，得到带入射源的方程
\[
 [E-H_{\rm eff}^{(+)}(E)]P|\Psi\rangle=H_{PQ}|\chi_E\rangle,
\]
其中

\begin{equation}
H_{\rm eff}^{(+)}(E)
=H_{PP}
+H_{PQ}G_Q^{(+)}(E)H_{QP}.
\label{eq:feshbach-effective}
\end{equation}

第二项

\[
\Sigma(E)
\equiv H_{PQ}G_Q^{(+)}(E)H_{QP}
\]

就是能量依赖自能。它的实部移动共振位置，连续谱耦合产生的虚部决定有限寿命。$H_{\rm eff}^{(+)}$ 的非厄米部分来自出射边界条件；完整哈密顿量仍自伴。寻找束缚态或经解析延拓后的共振极点时取零入射源；计算被连续波激发的散射态时，则须保留 $|\chi_E\rangle$。

\section{Fano 共振}
取

\[
P=|\phi\rangle\langle\phi|,
\qquad
H_0|\phi\rangle=E_\phi|\phi\rangle,
\]

并令连续态满足能量归一化

\[
\langle E|E'\rangle=\delta(E-E').
\]

定义耦合

\[
V_E\equiv\langle E|V|\phi\rangle.
\]

则标量自能为

\[
\Sigma(E)
=\int\dd E'\,
\frac{|V_{E'}|^2}{E-E'+\ii0^+}.
\]

Sokhotski--Plemelj 公式

\[
\frac1{x+\ii0^+}
=\mathcal P\frac1x-\ii\pi\delta(x)
\]

给出

\begin{equation}
\Sigma(E)
=\Delta(E)-\frac{\ii}{2}\Gamma(E),
\label{eq:self-energy-split}
\end{equation}

其中

\begin{equation}
\Delta(E)
=\mathcal P\int\dd E'\,
\frac{|V_{E'}|^2}{E-E'},
\qquad
\Gamma(E)=2\pi|V_E|^2.
\label{eq:shift-width}
\end{equation}

\(\Gamma\) 的具体前因子依赖连续态归一化；如果改用动量归一化，态密度 Jacobian 必须显式出现。

共振极点由

\[
E-E_\phi-\Sigma(E)=0
\]

决定。若 \(\Delta(E),\Gamma(E)\) 在共振附近缓慢变化，定义

\[
E_R=E_\phi+\Delta(E_R),
\]

就得到近似极点

\[
E_{\rm pole}\simeq E_R-\ii\Gamma(E_R)/2.
\]

对应时间振幅 \(\exp(-\ii E_{\rm pole}t/\hbar)\) 的概率按 \(\exp(-\Gamma t/\hbar)\) 衰减，所以寿命

\[
\tau\simeq\frac{\hbar}{\Gamma}.
\]

若自能的能量依赖不能忽略，极点附近传播子不是简单的单位留数。设 $z_p-E_\phi-\Sigma(z_p)=0$，则
\begin{equation}
G_P(z)\simeq\frac{Z_p}{z-z_p},
\qquad
Z_p=\frac{1}{1-\Sigma'(z_p)}.
\label{eq:resonance-pole-residue}
\end{equation}
$Z_p$ 衡量裸离散态在穿衣共振中的极点权重；宽共振或阈值附近 $\Sigma'(z)$ 较大时，把 $Z_p$ 近似为 1 会丢失重要谱权重信息。

复极点确定寿命，但探针同时可以直接激发连续态与间接激发离散态，两条路径的干涉还会改变线型。为分别计算态与探针振幅，先把模型的定义域和归一化固定下来。取
$\mathcal H=\mathbb C|\phi\rangle\oplus L^2(\mathcal I,\dd E)$，$\langle\phi|\phi\rangle=1$，连续态按 $\langle E|E'\rangle=\delta(E-E')$ 归一化。设 $V_E\in L^2(\mathcal I)$，则
\[
 V=\int_{\mathcal I}\dd E\,[V_E|E\rangle\langle\phi|+V_E^*|\phi\rangle\langle E|]
\]
是有界的厄米有限秩扰动，$H=H_0+V$ 在 $\Dom H_0$ 上自伴。连续态相位可逐能量选定，使 $V_E\ge0$ 在 $\mathcal I$ 上几乎处处成立；零耦合处的相位任意。这种可测相位变换保持能量归一化。以下讨论的能量位于连续谱内部，且 $V_E$ 在附近光滑、严格为正，所用主值边界值存在；因而主值积分中的 $V_{E'}$ 也使用同一个实耦合约定。

\begin{proposition}[单离散态与连续谱的散射态]
定义
\[
 x(E)=E-E_\phi-\Delta(E),\qquad
 \gamma(E)=\pi V_E^2=\Gamma(E)/2,\qquad
 \epsilon_E=x(E)/\gamma(E).
\]
具有单位入射连续振幅的散射态展开为
\begin{equation}
 |\Psi_E^{(\pm)}\rangle
 =a_E^{(\pm)}|\phi\rangle
 +\int_{\mathcal I}\dd E'\,b_E^{(\pm)}(E')|E'\rangle,
 \qquad a_E^{(\pm)}=\frac{V_E}{x(E)\pm\ii\gamma(E)},
 \label{eq:fano-state-expansion}
\end{equation}
其中
\begin{equation}
 b_E^{(\pm)}(E')
 =\delta(E-E')+
 \frac{a_E^{(\pm)}V_{E'}}{E-E'\pm\ii0}.
 \label{eq:fano-distribution-solution}
\end{equation}
将主值尾项合成
\[
 |\Phi_E\rangle=|\phi\rangle+
 \mathcal P\int_{\mathcal I}\dd E'\,\frac{V_{E'}}{E-E'}|E'\rangle,
\]
便得到同一散射态的表达
\begin{equation}
 |\Psi_E^{(\pm)}\rangle
 =\frac{\epsilon_E|E\rangle+(\pi V_E)^{-1}|\Phi_E\rangle}
        {\epsilon_E\pm\ii}.
 \label{eq:fano-normalized-state}
\end{equation}
这些是能量 $\delta$ 归一化的广义本征态，既不是两个可归一化态的有限线性组合，也不含探针参数。
\end{proposition}
\begin{proof}
把 $H$ 代入 $(E-H)|\Psi_E^{(\pm)}\rangle=0$，分别投影到离散态和连续态，得到
\[
 (E-E_\phi)a_E=\int\dd E'\,V_{E'}b_E(E'),
 \qquad (E-E')b_E(E')=V_{E'}a_E.
\]
后一个方程在分布意义下允许 $\delta(E-E')$ 齐次项。单位入射振幅将其系数固定为一，出射或入射边界则固定 $\pm\ii0$，给出\eqnrefs{eq:fano-distribution-solution}。代入离散方程，使用
$\Sigma(E\pm\ii0)=\Delta(E)\mp\ii\pi V_E^2$，求得 $a_E^{(\pm)}$。

为得到\eqnrefs{eq:fano-normalized-state}，用 Plemelj 公式拆开连续分量：$|E\rangle$ 的系数变成
\[
 1\mp\ii\pi V_Ea_E^{(\pm)}
 =\frac{x(E)}{x(E)\pm\ii\gamma(E)},
\]
而主值尾项与离散分量共同组成 $a_E^{(\pm)}|\Phi_E\rangle$。同时除以 $\gamma(E)$ 即得结果。

归一化使用有限秩散射的波算符结论：自伴有限秩扰动的波算符在 $H_0$ 的绝对连续子空间上存在并保持内积；作用于 $|E\rangle$ 因而保留 $\delta(E-E')$ 归一化。这是所引用的散射定理输入，上面的计算证明其在本模型中的显式形式。\footnote{有限秩波算符与绝对连续谱等价的定理陈述见 C. Liaw，\href{https://www.math.stonybrook.edu/~dimitriosnt/caft2018/slides/Liaw.pdf}{Finite Rank Perturbations}，第 12 页。}
离散分量还可由预解式独立核对：
\[
 -\frac1\pi\operatorname{Im}\frac1{x(E)+\ii\gamma(E)}
 =\frac{V_E^2}{x(E)^2+\gamma(E)^2}
 =|a_E^{(+)}|^2.
\]
这给出裸离散态在连续谱中的谱密度；若还存在真正束缚态，其离散权重须另外加入谱和。
\end{proof}

对于单通道弹性散射，由 $S=1-2\pi\ii V_Ea_E^{(+)}$ 得
\[
 S_{\rm res}(E)=\frac{x(E)-\ii\gamma(E)}{x(E)+\ii\gamma(E)}
 =\frac{\epsilon_E-\ii}{\epsilon_E+\ii}.
\]
实能量上 $|S_{\rm res}|=1$。设 $S_{\rm res}=\ee^{2\ii\delta_E}$，则
\begin{equation}
 \tan\delta_E=-1/\epsilon_E,
 \label{eq:fano-phase-shift}
\end{equation}
分支按跨过共振时连续增加的相移选定。

\begin{proposition}[探针振幅与 Fano 线型]
设 $T$ 为跃迁算符，$|i\rangle$ 为初态，$M_{\rm bg}(E)=\langle E|T|i\rangle\ne0$。定义
\[
 \widetilde M_\phi(E)=\langle\Phi_E|T|i\rangle,
 \qquad q(E)=\frac{\widetilde M_\phi(E)}{\pi V_EM_{\rm bg}(E)}.
\]
则以入射归一化的最终散射态计算跃迁，得到
\[
 \frac{\langle\Psi_E^{(-)}|T|i\rangle}{M_{\rm bg}(E)}
 =\frac{\epsilon_E+q(E)}{\epsilon_E+\ii}.
\]
在单通道且 $q$ 为实数的条件下，截面为
\begin{equation}
 \sigma(E)=\sigma_{\rm bg}(E)
 \frac{[q(E)+\epsilon_E]^2}{1+\epsilon_E^2}.
 \label{eq:fano-profile}
\end{equation}
复数 $q$ 时，保留振幅的模平方：
\begin{equation}
\sigma(E)=\sigma_{\rm bg}(E)\frac{|q(E)+\epsilon_E|^2}{1+\epsilon_E^2}.
\label{eq:phys-rl-fano-profile}
\end{equation}
\end{proposition}
\begin{proof}
对\eqnrefs{eq:fano-normalized-state}的负号支取伴随，得到
\[
 \langle\Psi_E^{(-)}|T|i\rangle
 =\frac{\epsilon_EM_{\rm bg}+(\pi V_E)^{-1}\widetilde M_\phi}
 {\epsilon_E+\ii}.
\]
提出 $M_{\rm bg}$ 后即得振幅比。若改用实驻波归一化，两个分量同时乘一个相位，截面仍相同。$q$ 由探针的两个矩阵元之比决定，所以改变探针可以改变线型，却不改变 $H$ 的散射态和共振极点。若 $M_{\rm bg}=0$，应直接使用上面的未约化振幅，$q$ 的比值定义不再适用。
\end{proof}

实验中常用常数 $q,\Gamma,E_R$ 拟合窄共振。这个近似要求 $V_E$、背景振幅和 $\Delta(E)$ 在一个宽度内变化很小。由
$x(E)=E-E_R-[\Delta(E)-\Delta(E_R)]$
可见，$|\Delta'(E_R)|\ll1$ 时才有 $\epsilon_E\simeq(E-E_R)/(\Gamma(E_R)/2)$；阈值附近应保留能量依赖。单通道时间反演对称、探针矩阵元可同时取实的条件给出实 $q$，存在不可消去的相位时则保留复 $q$。

\begin{workedexample}[相同共振极点下的两种探针线型]
取缓变背景与 $\Gamma>0$，分别令 $q=0$、$q=1$。求截面零点、共振中心值及 $q=1$ 的峰值位置。

\textbf{解。}$q=0$ 时 $\sigma/\sigma_{\rm bg}=\epsilon^2/(1+\epsilon^2)$，在 $E=E_R$ 完全相消。$q=1$ 时零点为 $\epsilon=-1$，即 $E=E_R-\Gamma/2$；中心 $\epsilon=0$ 的截面比为一。对线型求导，
\[
 \frac{\dd}{\dd\epsilon}\frac{(1+\epsilon)^2}{1+\epsilon^2}
 =\frac{2(1+\epsilon)(1-\epsilon)}{(1+\epsilon^2)^2},
\]
所以另一极值在 $\epsilon=1$，峰值为二。两种探针的极点都为 $E_R-\ii\Gamma/2$，寿命都为 $\hbar/\Gamma$；变化的是干涉零点及峰值，而不是系统的衰变宽度。
\end{workedexample}

不同 \(q\) 的线型见\figref{fig:fano-lineshape}。

\begin{figure}[tbp]
\centering
\includegraphics[width=.83\textwidth]{fano_lineshape.pdf}
\caption{不同 Fano 参数 \(q\) 对归一化线型的影响。\(|q|\gg1\) 时共振趋向 Lorentz 峰，\(q=0\) 时共振中心成为反共振零点。}
\label{fig:fano-lineshape}
\end{figure}

冷原子 Feshbach 共振使用的是同一投影结构，只是空间含义对调：\(P\) 是开放散射通道，\(Q\) 包含闭通道束缚态。若闭通道中只有一个束缚态 \(|\phi_c\rangle\) 靠近开放通道阈值，则开通道的有效相互作用包含

\[
V_{\rm eff}(E,B)
\supset
\frac{W|\phi_c\rangle\langle\phi_c|W^\dagger}
{E-E_c(B)-\Sigma_c(E)}.
\]

在超低能 s 波区，散射振幅满足有效程展开

\[
k\cot\delta_0(k)
=-\frac1a+\frac12r_ek^2+\cdots.
\]

若闭通道束缚态相对于阈值的失谐在窄磁场区间近似线性，

\[
E_c(B)-E_{\rm th}
\simeq\delta\mu(B-B_c),
\]

并把实自能产生的阈值位移吸收到新的共振位置 \(B_0\)，则零能 \(T\)-矩阵中的共振项必然产生一个 \((B-B_0)^{-1}\) 极点。因此散射长度可以参数化为

\begin{equation}
a(B)
=a_{\rm bg}
\left(1-\frac{\Delta B}{B-B_0}\right).
\label{eq:feshbach-scattering-length}
\end{equation}

\(a_{\rm bg}\) 吸收所有缓慢变化的开放通道背景，\(\Delta B\) 把耦合强度和磁矩差压缩成实验可测的磁场宽度。\eqnrefs{eq:feshbach-scattering-length} 是“单闭通道束缚态 + 低能 s 波 + 背景缓慢变化”的阈值近似，而不是一般 Feshbach 理论的定义。

\begin{proof}
磁场通过磁矩差 $\delta\mu=\mu_{\rm closed}-\mu_{\rm open}$ 移动闭通道束缚态，$E_c(B)\approx\delta\mu(B-B_c)$。零能附近单闭通道消元后，共振部分为
\begin{equation*}
 f_{\rm res}(0,B)
 =\frac{C_f}
 {-\delta\mu(B-B_c)-\Delta_c(0)},
\end{equation*}
其中 $C_f$ 汇总耦合常数，$\Delta_c(0)$ 为实自能位移。极点处
\begin{equation*}
 B_0
 =B_c-\frac{\Delta_c(0)}{\delta\mu},
\end{equation*}
于是散射振幅分解为背景加共振极点：
\begin{equation*}
 f(0,B)
 =f_{\rm bg}(0)-\frac{C_f/\delta\mu}{B-B_0}.
\end{equation*}


三维 s 波阈值 $f_0(k\to0)=-a$，于是
\begin{equation}
 a(B)=a_{\rm bg}-\frac{C_a}{B-B_0},
 \qquad
 C_a\equiv-\frac{C_f}{\delta\mu}.
 \label{eq:feshbach-scattering-length-additive}
\end{equation}


只要 $a_{\rm bg}\neq0$，定义 $\Delta B=C_a/a_{\rm bg}$，则
\begin{equation*}
 a(B)=a_{\rm bg}\left(1-\frac{\Delta B}{B-B_0}\right),
\end{equation*}
即\eqnrefs{eq:feshbach-scattering-length}。$a(B)$ 在 $B=B_0$ 处发散、在 $B=B_0+\Delta B$ 处过零。

\end{proof}

\section{Landau--Zener 跃迁}

\begin{equation}
H(t)
=\begin{pmatrix}
s_{\rm LZ} t/2&V\\
V^*&-s_{\rm LZ} t/2
\end{pmatrix},
\label{eq:lz-hamiltonian}
\end{equation}

其中 \(s_{\rm LZ}\) 的量纲是能量/时间，\(2|V|\) 是避免交叉的最小绝热能隙。瞬时绝热能量

\[
E_\pm(t)
=\pm\sqrt{(s_{\rm LZ} t/2)^2+|V|^2}.
\]

\figref{fig:landau-zener} 展示了非绝热交叉被恒定耦合打开的过程。

\begin{figure}[tbp]
\centering
\includegraphics[width=.83\textwidth]{landau_zener.pdf}
\caption{Landau--Zener 模型的 diabatic 与 adiabatic 能级。常数耦合 \(V\) 把线性交叉打开为最小间隙 \(2|V|\)。}
\label{fig:landau-zener}
\end{figure}

通过常数相位变换可令 \(V>0\) 为实数。两分量方程为

\begin{align*}
\ii\hbar\dot c_1
&=\frac{s_{\rm LZ} t}{2}c_1+Vc_2,
\\
\ii\hbar\dot c_2
&=Vc_1-\frac{s_{\rm LZ} t}{2}c_2.
\end{align*}

消去 \(c_2\)。从第一式

\[
c_2
=\frac{\ii\hbar\dot c_1-(s_{\rm LZ} t/2)c_1}{V}.
\]

对第一式再求导并代入第二式，\(\dot c_1\) 项恰好相消，得到

\begin{equation}
\ddot c_1
+\left[
\frac{s_{\rm LZ}^2t^2}{4\hbar^2}
+\frac{V^2}{\hbar^2}
+\frac{\ii\,s_{\rm LZ}}{2\hbar}
\right]c_1=0.
\label{eq:lz-second-order}
\end{equation}

对 \(s_{\rm LZ}>0\) 定义

\[
z=\ee^{-\ii\pi/4}
\sqrt{\frac{s_{\rm LZ}}{\hbar}}t,
\qquad
\delta\equiv\frac{V^2}{\hbar s_{\rm LZ}}.
\]

由于

\[
\frac{\dd^2}{\dd t^2}
=-\ii\frac{s_{\rm LZ}}{\hbar}
\frac{\dd^2}{\dd z^2},
\qquad
 t^2=\ii\frac{\hbar}{s_{\rm LZ}}z^2,
\]

\eqnrefs{eq:lz-second-order} 变成

\begin{equation}
\frac{\dd^2c_1}{\dd z^2}
+\left(
\nu+\frac12-\frac{z^2}{4}
\right)c_1=0,
\qquad
\nu=-1+\ii\delta.
\label{eq:lz-weber}
\end{equation}

这就是抛物柱 (Weber) 方程，其独立解可写为 \(D_\nu(z)\) 与适当旋转后的 \(D_\nu(-z)\)。

\section{Stokes 连接、跃迁指数与 Stückelberg 相位}

Landau--Zener 概率的非平凡之处不在“解是特殊函数”，而在如何把 \(t\to-\infty\) 的入射边界条件穿过复平面的 Stokes 扇区连接到 \(t\to+\infty\)。利用抛物柱函数的标准连接公式，保持非绝热特征标的概率为

\begin{equation}
P_{\rm LZ}
=\ee^{-2\pi\delta}
=\exp\left[-\frac{2\pi|V|^2}{\hbar|s_{\rm LZ}|}\right].
\label{eq:landau-zener}
\end{equation}

\begin{proof}
\eqnrefs{eq:lz-weber} 的解 $D_\nu(z)$（$\nu=-1+\ii\delta$）在不同复 $z$ 扇区有不同的渐近行为。Stokes 现象是指：同一解在不同扇区的渐近展开系数不同，连接公式描述跨扇区时这些系数如何跳变。在前面固定的 $z=\ee^{-\ii\pi/4}\sqrt{s_{\rm LZ}/\hbar}\,t$ 约定下，$t\to-\infty$ 位于 $\arg z=3\pi/4$，$t\to+\infty$ 位于 $\arg z=-\pi/4$。两端位于不同扇区，须使用相应的抛物柱函数连接公式。

入射/出射解按单位概率流归一化后，连接矩阵对角振幅模为
\[
 |S_{11}|=\ee^{-\pi\delta},
\]
$\ee^{-\pi\delta}$ 的指数结构来自抛物柱函数 $D_\nu$ 在跨扇区时的次主导指数项系数，是 Stokes 现象的精确体现。

Stokes 乘子可写成
\[
 C_\delta
 =\frac{\sqrt{2\pi}\,\ee^{-\pi\delta/2}}
 {\sqrt{\delta}\,\Gamma(-\ii\delta)}
 \times \ee^{\ii\varphi_\delta},
\]
其中 \(\varphi_\delta\) 只影响相位。$\Gamma$ 函数 $|\Gamma(-\ii\delta)|^2$ 的模由两个恒等式决定：递推 $\Gamma(1-\ii\delta)=(-\ii\delta)\Gamma(-\ii\delta)$ 和反射公式给出的
\[
 \Gamma(1-\ii\delta)=(-\ii\delta)\Gamma(-\ii\delta),
 \qquad
 |\Gamma(1-\ii\delta)|^2
 =\frac{\pi\delta}{\sinh(\pi\delta)},
\]
后者可由反射公式直接核对：
\[
 \Gamma(\ii\delta)\Gamma(1-\ii\delta)
 =\frac{\pi}{\sin(\pi\ii\delta)}
 =-\frac{\ii\pi}{\sinh(\pi\delta)}.
\]
再用 $\Gamma(1-\ii\delta)=-\ii\delta\Gamma(-\ii\delta)$ 和共轭关系，得到
$|\Gamma(-\ii\delta)|^2=\pi/[\delta\sinh(\pi\delta)]$，乘上 $\delta^2$ 即得所用的 $|\Gamma(1-\ii\delta)|^2$。

代入 $|C_\delta|^2$，
\begin{align*}
 |C_\delta|^2
 &=\frac{2\pi\ee^{-\pi\delta}}
 {\delta|\Gamma(-\ii\delta)|^2}\\
 &=2\ee^{-\pi\delta}\sinh(\pi\delta)\\
 &=1-\ee^{-2\pi\delta}.
\end{align*}
最后一步用 $\sinh(\pi\delta)=(\ee^{\pi\delta}-\ee^{-\pi\delta})/2$。因此单位流守恒给出
\[
 |S_{11}|^2=\ee^{-2\pi\delta},
 \qquad
 |S_{21}|^2=1-\ee^{-2\pi\delta}.
\]
前者即 \eqnrefs{eq:landau-zener}。$|S_{11}|^2+|S_{21}|^2=1$ 验证幺正性。连接系数的相位整理后得 Stokes 相位 \(\phi_S\)（\eqnrefs{eq:lz-stokes-phase}），在 Stückelberg 干涉中与动力学相位叠加。$\delta=V^2/(\hbar s_{\rm LZ})$ 使 $s_{\rm LZ}\to0$ 时 $P_{\rm LZ}\to0$、$s_{\rm LZ}\to\infty$ 时 $P_{\rm LZ}\to1$。
\end{proof}

控制参数

\[
\delta=\frac{|V|^2}{\hbar|s_{\rm LZ}|}
\]

比较“能隙平方”与“扫过速率”。\(\delta\gg1\) 时交叉被缓慢穿过，系统绝热跟随；\(\delta\ll1\) 时扫过太快，状态来不及混合，几乎保持原非绝热特征标。

如果不仅关心单次跃迁概率，还关心多次交叉的 Stückelberg 干涉，就必须保留振幅相位。单次 Landau--Zener 散射矩阵除了 \(P_{\rm LZ}\) 外还包含 Stokes 相

\begin{equation}
\phi_S
=\frac\pi4
+\arg\Gamma(1-\ii\delta)
+\delta(\ln\delta-1).
\label{eq:lz-stokes-phase}
\end{equation}

在一种常见相位约定下可写

\[
S_{\rm LZ}
=\begin{pmatrix}
\sqrt{P_{\rm LZ}}\ee^{-\ii\phi_S}&-\sqrt{1-P_{\rm LZ}}\\
\sqrt{1-P_{\rm LZ}}&\sqrt{P_{\rm LZ}}\ee^{+\ii\phi_S}
\end{pmatrix}.
\]

两次避免交叉之间积累的动力学相位与 \(\phi_S\) 干涉，就形成 St\"uckelberg 干涉条纹。

\begin{note}[Landau--Zener 公式何时不能直接用]
\eqnrefs{eq:landau-zener} 要求：非绝热能差在相关时间窗内可线性化、耦合 \(V\) 近似常数、演化从远离交叉的 \(t\to-\infty\) 延伸到 \(t\to+\infty\)，并且系统在这段时间内可视为封闭二能级系统。有限脉冲、非线性扫频、多个近邻能级或显著耗散都会改变连接问题，应回到一般含时传播子或多能级模型。
\end{note}

\section{Feshbach 共振与 Landau--Zener 跃迁的实验数值}\label{sec:resonance-experiments}

Feshbach 共振在超冷原子物理中用于精确调控散射长度 \(a(B)\)：
\[
a(B)=a_{\rm bg}\left(1-\frac{\Delta B}{B-B_0}\right)
\]
其中 \(a_{\rm bg}\) 是背景散射长度，\(B_0\) 是共振磁场，\(\Delta B\) 是带符号的宽度参数；其正负由这一公式的约定决定。表中的长度以 Bohr 半径 \(a_0\) 为单位。锂的参数取 2005 年标定的近似值，适于比较共振两侧和数量级；精密实验需使用相应碰撞通道的更新标定。\footnote{M.~Bartenstein 等，\href{https://arxiv.org/abs/cond-mat/0408673}{射频谱测定的锂冷碰撞参数}；原文使用 \(1+\Delta/(B-B_0)\)，其正宽度 \(300\,\mathrm G\) 对应此处的 \(\Delta B=-300\,\mathrm G\)。}

\begin{center}
\begin{tabular}{lcccc}
\hline
体系 & \(B_0\) (G) & \(\Delta B\) (G) & \(a_{\rm bg}/a_0\) & 实验组 \\
\hline
\(^6\)Li \(|1\rangle+|2\rangle\) & 834.1 & \(-300\) & \(-1405\) & Innsbruck (2005) \\
\(^{40}\)K \(|9/2,-9/2\rangle\) & 202.1 & 7.8 & 174 & JILA (2004) \\
\(^{87}\)Rb \(|1,1\rangle+|1,0\rangle\) & 1007.4 & 0.21 & 99 & MIT (2004) \\
\(^{23}\)Na \(|1,1\rangle\) & 907.0 & 1.0 & 52 & MIT (2003) \\
\hline
\end{tabular}
\end{center}

\begin{exbox}[\(^6\)Li 的宽共振与 BEC--BCS 渡越]
取上表的 $B_0=834.1\,\mathrm G$、$\Delta B=-300\,\mathrm G$、$a_{\rm bg}=-1405a_0$。求共振附近 $B=B_0\mp10\,\mathrm G$ 的散射长度，判断哪一侧支持浅束缚分子。
\textbf{解。} 直接代入得
\[
 a(B_0-10\,\mathrm G)=(-1405a_0)(1-30)=40745a_0>0,
\]
\[
 a(B_0+10\,\mathrm G)=(-1405a_0)(1+30)=-43555a_0<0.
\]
因而共振附近的低场侧 $B<B_0$ 对应 $a>0$，支持浅束缚 Feshbach 分子，并可进入分子凝聚的一侧；高场侧 $B>B_0$ 对应 $a<0$，是 Cooper 配对的 BCS 一侧。靠近阈值的分子结合能为 $E_b\simeq\hbar^2/(ma^2)$，这些分子是弱束缚的，并非深结合分子。
当 $B\to B_0$ 时，$|a|\to\infty$，低能二体 $s$ 波截面趋于 $4\pi/k^2$。在有效程和相互作用范围可忽略的稀薄、宽共振极限中，散射长度不再引入有限长度尺度；多体问题仍具有密度确定的 $k_F^{-1}$ 和热波长等尺度，热力学量可写为费米能乘以 $T/T_F$ 的无量纲函数。这个标度结论不能解释成密度和温度本身消失。
\end{exbox}

\begin{center}
\begin{tabular}{lccc}
\hline
体系 & \(\delta_{\rm LZ}=V^2/(\hbar|\dot\Delta|)\) & 跃迁概率 & 实验 \\
\hline
超导量子比特 (Rabi) & 0.1--10 & 0--1 可调 & Nakamura (2003) \\
Rydberg 原子 Stark & 0.5--5 & 0.1--0.9 & Gallagher (1980s) \\
分子束共振 & \(\sim1\) & 0.5 附近 & 分子物理 \\
\hline
\end{tabular}
\end{center}

\begin{exbox}[给定量子比特能隙与扫速的跃迁概率]
取二能级最小能隙 $2V/h=100\,\mathrm{MHz}$，透热能差扫描速率 $|\dot\Delta|/h=1\,\mathrm{GHz/ns}$。计算留在原透热态的概率，并判断该扫速接近哪一极限。

\textbf{解。}先统一频率与时间单位：$V/h=5.0\times10^7\,\mathrm{s^{-1}}$，$|\dot\Delta|/h=10^{18}\,\mathrm{s^{-2}}$。使用 $h/\hbar=2\pi$，得到
\[
 \delta_{\rm LZ}=2\pi\frac{(V/h)^2}{|\dot\Delta|/h}
 =0.01571,\qquad P_{\rm LZ}=\ee^{-2\pi\delta_{\rm LZ}}\simeq0.9060.
\]
因而约 $90.6\%$ 留在原透热态，约 $9.4\%$ 沿绝热分支转移，这组参数靠近快扫极限。将扫速再提高一百倍，$\delta_{\rm LZ}=1.571\times10^{-4}$、$P_{\rm LZ}\simeq0.9990$。要得到强绝热跟随，需降低扫速，使 $V^2/(\hbar|\dot\Delta|)\gg1$。
\end{exbox}

Fano 共振的散射截面
\[
\sigma(\epsilon)=\sigma_0\frac{(q+\epsilon)^2}{1+\epsilon^2},\qquad\epsilon=\frac{E-E_r}{\Gamma/2}
\]
其中 \(q\) 是 Fano 不对称参数，\(\Gamma\) 是共振宽度。

\begin{center}
\begin{tabular}{lcc}
\hline
体系 & \(q\) & \(\Gamma\) \\
\hline
He 原子 \(2s2p\,{}^1P\) 自电离 & \(-2.8\) & 0.012 eV \\
量子点 Kondo & 0.5 & \(k_BT_K\) \\
等离子体 Fano & 1--10 & 0.1--1 eV \\
\hline
\end{tabular}
\end{center}

\begin{exbox}[He 自电离共振的 Fano 谱]
取负 $q$ 共振的线型模型 $q=-2.8$、$\Gamma=0.012\,\mathrm{eV}$，背景在共振附近为常数。求谱零点、峰值相对 $E_r$ 的位置、峰值与背景的比值，以及该宽度对应的寿命。

\textbf{解。}令 $F(\epsilon)=(q+\epsilon)^2/(1+\epsilon^2)$，则
\[
 F'(\epsilon)=\frac{2(q+\epsilon)(1-q\epsilon)}{(1+\epsilon^2)^2}.
\]
零点为 $\epsilon_0=-q=2.8$，其能量在共振中心上方：
\[
 E_0-E_r=\frac\Gamma2\epsilon_0=0.0168\,\mathrm{eV}.
\]
另一驻点为 $\epsilon_{\rm peak}=1/q=-0.3571$，所以
\[
 E_{\rm peak}-E_r=\frac\Gamma{2q}\simeq-0.00214\,\mathrm{eV},
 \qquad F(\epsilon_{\rm peak})=1+q^2=8.84.
\]
由于 $F\to1$ 在远离共振时恢复背景，这个非零驻点是峰值。最后由极点 $E_r-\ii\Gamma/2$ 得
\[
 \tau=\hbar/\Gamma
 =\frac{6.5821\times10^{-16}\,\mathrm{eV\,s}}{0.012\,\mathrm{eV}}
 \simeq54.9\,\mathrm{fs}.
\]
峰位、零点和极点的实部分别描述线型最大值、干涉相消和共振中心；同一宽度决定寿命，但峰位与零点还随探针参数 $q$ 改变。
\end{exbox}

\begin{proposition}[冻结宽度下的共振相位与延迟]
在单通道弹性散射中，若一个共振宽度内可冻结 $V_E$ 和实自能，令 $\epsilon=(E-E_r)/(\Gamma/2)$，则
\[
 S(E)=\ee^{2\ii\delta_{\rm bg}(E)}\frac{\epsilon-\ii}{\epsilon+\ii},
 \qquad
 \tau_{\rm res}(E)=\frac{4\hbar}{\Gamma(1+\epsilon^2)}.
\]
这里 $\tau_{\rm res}$ 为共振对 Wigner--Smith 延迟的贡献，总延迟另加背景项 $2\hbar\delta_{\rm bg}'(E)$。两者均须由系统的散射矩阵计算，探针参数 $q$ 不进入共振因子。
\end{proposition}

\begin{proof}[证明（共振相移与时间延迟）]

先由能量归一化散射态求出 $S$ 矩阵，再对相移求导得到延迟。在共振宽度内冻结 $V_E$ 与实自能，将\eqnrefs{eq:fano-normalized-state}的出射支写为
\begin{equation}
|\Psi_E^{(+)}\rangle
=\frac{\epsilon|E\rangle+(\pi V_E)^{-1}|\Phi_E\rangle}{\epsilon+\ii},
\qquad
\epsilon=\frac{E-E_r}{\Gamma/2},
\label{eq:fano-state-s-matrix}
\end{equation}
此处 $|\Phi_E\rangle$ 包含完整主值连续谱尾项。由 $S_{\rm res}=1-2\pi\ii V_Ea_E^{(+)}$，并乘上缓变的背景相位，得到
\begin{equation}
S(\epsilon)=\ee^{2\ii\delta_{\rm bg}}\,\frac{\epsilon-\ii}{\epsilon+\ii},
\label{eq:fano-s-matrix}
\end{equation}
式 \eqnrefs{eq:fano-s-matrix} 的极点位于 \(\epsilon=-\ii\)，即 \(E=E_r-\ii\Gamma/2\)。\(q\) 进入具体跃迁算符的 Fano 线型，而不是这一极点位置。

由 \(S=\ee^{2\ii\delta}\) 提取 \(\delta=\delta_{\rm bg}+\delta_{\rm res}\)，共振部分满足
\begin{equation}
\delta_{\rm res}(\epsilon)=-\arctan\!\frac{1}{\epsilon},
\label{eq:fano-resonance-phase-shift}
\end{equation}
其中反正切的分支取成跨过 \(\epsilon=0\) 时相移连续增加 \(\pi\)。因此相移中心由 \(E_r\) 决定，与 Fano 零点 \(\epsilon=-q\) 一般不重合。

Wigner--Smith 时间延迟为
\begin{equation}
\tau_{\rm res}(\epsilon)=2\hbar\,\frac{\dd\delta_{\rm res}}{\dd E}
=\frac{4\hbar}{\Gamma}\,\frac{\dd\delta_{\rm res}}{\dd\epsilon},
\label{eq:wigner-smith-definition}
\end{equation}
因子 \(4\hbar/\Gamma\) 来自 \(\dd E=(\Gamma/2)\dd\epsilon\)。由\eqnrefs{eq:fano-resonance-phase-shift}，
\begin{equation}
\frac{\dd\delta_{\rm res}}{\dd\epsilon}
=\frac{1}{1+\epsilon^2},
\end{equation}
故
\begin{equation}
\tau_{\rm res}(\epsilon)=\frac{4\hbar}{\Gamma}\,\frac{1}{1+\epsilon^2}.
\label{eq:fano-wigner-smith-delay}
\end{equation}
共振贡献在 $\epsilon=0$ 取最大值 $4\hbar/\Gamma$；总延迟还须加上 $2\hbar\,\dd\delta_{\rm bg}/\dd E$。探针的 Fano 零点由 $q$ 决定，两者仅在 $q=0$ 时重合。背景导数为负时，总延迟可含负修正；阈值或宽度随能量显著变化时还须对完整 $S(E)$ 求导。

\end{proof}

\subsection{多离散态的宽度矩阵与暗态}

单态模型的宽度是一个非负数。若保留多个离散态 $|\phi_n\rangle$，同一个连续通道可能同时耦合它们，衰变就由矩阵控制。设连续态按
$\langle E,\alpha|E',\beta\rangle=\delta_{\alpha\beta}\delta(E-E')$
归一化，并记 $W_{\alpha n}(E)=\langle E,\alpha|V|\phi_n\rangle$。分块消元给出
\[
 \Sigma_{nm}(E+\ii0)
 =\sum_\alpha\int\dd E'\,
 \frac{W_{\alpha n}(E')^*W_{\alpha m}(E')}{E-E'+\ii0}
 =\Delta_{nm}(E)-\frac\ii2\Gamma_{nm}(E).
\]

\begin{proposition}[连续谱产生的宽度矩阵]
在上述能量归一化下，
\[
 \Gamma(E)=2\pi W(E)^\dagger W(E)\succeq0.
\]
散射极点由 $\det[z-H_{PP}-\Sigma(z)]=0$ 决定；$\Gamma$ 的零空间只给出零直接衰变方向。若 $v\in\ker W(E)$ 同时受实有效哈密顿量 $H_{PP}+\Delta(E)$ 保持，才可能构成该能量处的稳定暗态。
\end{proposition}
\begin{proof}
对每个矩阵元使用 Plemelj 公式，虚部为 $-\pi\sum_\alpha W_{\alpha n}^*W_{\alpha m}$，故得到 $\Gamma$。对任意系数向量 $v$，
\[
 v^\dagger\Gamma v=2\pi\sum_\alpha
 \left|\sum_nW_{\alpha n}v_n\right|^2\ge0.
\]
若 $Wv=0$，该线性组合的直接连续谱耦合相消；若实有效哈密顿量再将它混到亮态中，它仍可间接衰变。因此零宽度方向与真正稳定的本征态须同时检查。多通道探针跃迁还需各通道背景与离散路径的振幅；不可只给每条通道一条互不干涉的标量 Fano 曲线。
\end{proof}

\begin{workedexample}[两条简并离散态耦合到同一连续通道]
取 $H_{PP}=E_\phi I_2$，并令两态对所有连续谱能量都有相同实耦合 $W(E)=(v_E,v_E)$。求亮态、暗态与它们的衰变宽度。

\textbf{解。}宽度和实自能都正比于全一矩阵：
\[
 \Gamma(E)=2\pi v_E^2\begin{pmatrix}1&1\\1&1\end{pmatrix},\qquad
 \Delta(E)=\mathcal P\int\frac{v_{E'}^2}{E-E'}\dd E'
                  \begin{pmatrix}1&1\\1&1\end{pmatrix}.
\]
共同本征向量为 $|B\rangle=(|\phi_1\rangle+|\phi_2\rangle)/\sqrt2$ 与
$|D\rangle=(|\phi_1\rangle-|\phi_2\rangle)/\sqrt2$。前者耦合为 $\sqrt2v_E$，宽度为 $\Gamma_B=4\pi v_E^2$；后者对所有连续能量的耦合均为零，故 $\Gamma_D=0$、$\Delta_D=0$。

$H_{PP}$ 也保持这两条方向，所以 $H|D\rangle=E_\phi|D\rangle$：即使 $E_\phi$ 位于连续谱中，它仍是可归一化的暗本征态。单条离散态的宽度为 $2\pi v_E^2$，亮态宽度正好加倍。若两个裸能量失去简并，$H_{PP}$ 在 $B,D$ 基中产生非对角元，暗态通常获得间接衰变通道。
\end{workedexample}
